%%
%% part1-getting-started/ch03-first-book.tex
%%
%% Chapter 3: Your First Book — Build a minimal but complete
%% book with MatinBook, from cover to back cover.
%%

\chapter{اولین کتاب}
\label{chap:first-book}

\section{چرا این موضوع مهم است؟}
\label{sec:first-book-why}

پس از نصب موفق \lr{MatinBook}، وقت آن رسیده که
اولین کتاب خود را بسازید. هدف این فصل، نه آموزش
تمام قابلیت‌های \lr{MatinBook}، بلکه آشنایی با
\textbf{ساختار پایه‌ی یک کتاب} است. پس از خواندن
این فصل، شما یک کتاب کامل با جلد، پیشگفتار، فصل،
منابع و نمایه خواهید داشت.

این کتاب کوچک، پایه‌ای خواهد بود که در فصل‌های
بعدی، با افزودن قابلیت‌های پیشرفته‌تر، آن را
گسترش خواهید داد.

\mnote{برای شروع، نیازی به دانش عمیق \lr{LaTeX} نیست. کافی است دستورات این فصل را گام‌به‌گام دنبال کنید.}

\section{ساختار یک کتاب کامل}
\label{sec:first-book-structure}

یک کتاب حرفه‌ای در \lr{MatinBook} از پنج بخش
اصلی تشکیل شده است:

\begin{enumerate}
    \item \textbf{جلد جلو:} با دستور \cmd{makecover}.
    \item \textbf{پیش‌متن:} شامل شناسنامه، تقدیم،
    پیشگفتار، فهرست مطالب.
    \item \textbf{متن اصلی:} شامل فصل‌ها.
    \item \textbf{پس‌متن:} شامل منابع و نمایه.
    \item \textbf{جلد عقب:} با دستور \cmd{makebackcover}.
\end{enumerate}

\mnote{ترتیب این پنج بخش مهم است. اگر جابه‌جا شوند، کتاب به‌درستی کامپایل نمی‌شود.}

\section{ساخت فایل اصلی}
\label{sec:first-book-main}

یک فایل جدید به نام \file{my-first-book.tex} بسازید
و محتوای زیر را در آن قرار دهید:

\begin{latin}
\begin{minted}{latex}
\documentclass{matinbook}

\title{|\pc{اولین کتاب من}|}
\author{|\pc{نام من}|}
\date{|\pc{مهر ۱۴۰۵}|}

\booktitle{|\pc{اولین کتاب من}|}

\begin{document}

% |\pc{جلد جلو}|
\makecover
    {|\pc{اولین کتاب من}|}
    {|\pc{یک کتاب آزمایشی}|}
    {|\pc{نام من}|}
    {|\pc{مهر ۱۴۰۵}|}

% |\pc{پیش‌متن}|
\frontmatter
\frontmattergeometry

\maketitle
\tableofcontents

% |\pc{متن اصلی}|
\mainmatter
\mainmattergeometry

\chapter{|\pc{فصل اول}|}
\label{chap:first}

|\pc{این یک متن ساده است.}|

% |\pc{پس‌متن}|
\backmatter
\frontmattergeometry

% |\pc{جلد عقب}|
\makebackcover{|\pc{توضیحات پشت جلد...}|}

\end{document}
\end{minted}
\end{latin}

\mnote{نام فایل را می‌توانید هر چه می‌خواهید بگذارید. در این کتاب از \lr{my-first-book.tex} استفاده می‌کنیم.}

\section{کامپایل اولین کتاب}
\label{sec:first-book-compile}

برای ساخت فایل \lr{PDF}، از دستور زیر استفاده
کنید:

\begin{latin}
\begin{minted}{bash}
xelatex -shell-escape my-first-book.tex
xelatex -shell-escape my-first-book.tex
\end{minted}
\end{latin}

\mnote{دو بار کامپایل لازم است چون جلد با \lr{TikZ} ساخته می‌شود و به دو مرحله نیاز دارد.}

پس از اجرای این دستورات، یک فایل
\file{my-first-book.pdf} ساخته می‌شود. آن را باز
کنید و نتیجه را ببینید.

\mnote{اگر خطایی گرفتید، به \cref{app:troubleshooting} مراجعه کنید.}

\section{آشنایی با ساختار}
\label{sec:first-book-tour}

حالا که اولین کتاب خود را ساختید، بگذارید
ساختار آن را بخش‌به‌بخش بررسی کنیم.

\subsection{فرمان \cmd{documentclass}}
\label{sec:first-book-documentclass}

اولین خط هر سند \lr{LaTeX}، دستور
\cmd{documentclass} است:

\begin{latin}
\begin{minted}{latex}
\documentclass{matinbook}
\end{minted}
\end{latin}

این دستور به \lr{LaTeX} می‌گوید که از کلاس
\lr{MatinBook} استفاده کند. می‌توانید آپشن‌هایی
مانند \opt{draft} یا \opt{vazirmatn} را هم اضافه
کنید:

\begin{latin}
\begin{minted}{latex}
\documentclass[draft, vazirmatn]{matinbook}
\end{minted}
\end{latin}

\mnote{آپشن‌ها با کاما از هم جدا می‌شوند. ترتیب آن‌ها مهم نیست.}

\subsection{فراداده کتاب}
\label{sec:first-book-metadata}

بعد از \cmd{documentclass}، فراداده‌ی کتاب را
تعریف می‌کنیم:

\begin{latin}
\begin{minted}{latex}
\title{|\pc{اولین کتاب من}|}
\author{|\pc{نام من}|}
\date{|\pc{مهر ۱۴۰۵}|}

\booktitle{|\pc{اولین کتاب من}|}
\end{minted}
\end{latin}

\begin{itemize}
    \item \cmd{title}: عنوان کتاب (برای صفحه‌ی
    عنوان).
    \item \cmd{author}: نام نویسنده.
    \item \cmd{date}: تاریخ.
    \item \cmd{booktitle}: عنوان کتاب برای
    سرصفحه‌ی صفحات زوج.
\end{itemize}

\mnote{فراداده‌ی \cmd{title} و آرگومان اول \cmd{makecover} باید یکسان باشند. برای جلوگیری از ناسازگاری، آن‌ها را هماهنگ نگه دارید.}

\subsection{جلد جلو}
\label{sec:first-book-frontcover}

جلد جلو با دستور \cmd{makecover} ساخته می‌شود:

\begin{latin}
\begin{minted}{latex}
\makecover
    {|\pc{اولین کتاب من}|}
    {|\pc{یک کتاب آزمایشی}|}
    {|\pc{نام من}|}
    {|\pc{مهر ۱۴۰۵}|}
\end{minted}
\end{latin}

این دستور چهار آرگومان می‌گیرد:

\begin{enumerate}
    \item عنوان کتاب.
    \item زیرعنوان.
    \item نام نویسنده.
    \item تاریخ یا نسخه.
\end{enumerate}

\subsection{پیش‌متن}
\label{sec:first-book-frontmatter}

پیش‌متن با دستور \cmd{frontmatter} شروع می‌شود:

\begin{latin}
\begin{minted}{latex}
\frontmatter
\frontmattergeometry
\end{minted}
\end{latin}

دستور \cmd{frontmatter} شماره‌گذاری صفحات را به
حروف ابجد فارسی (الف، ب، ج، ...) تغییر می‌دهد.
دستور \cmd{frontmattergeometry} نیز هندسه‌ی
صفحه را به حالت متقارن تغییر می‌دهد.

\mnote{در پیش‌متن، حاشیه‌ی بیرونی پهن وجود ندارد. این حالت برای فهرست مطالب و پیشگفتار مناسب‌تر است.}

\subsection{صفحه‌ی عنوان}
\label{sec:first-book-titlepage}

صفحه‌ی عنوان با دستور \cmd{maketitle} ساخته می‌شود:

\begin{latin}
\begin{minted}{latex}
\maketitle
\end{minted}
\end{latin}

این دستور از فراداده‌ی \cmd{title}، \cmd{author} و
\cmd{date} استفاده می‌کند.

\subsection{فهرست مطالب}
\label{sec:first-book-toc}

فهرست مطالب با دستور \cmd{tableofcontents} ساخته
می‌شود:

\begin{latin}
\begin{minted}{latex}
\tableofcontents
\end{minted}
\end{latin}

\mnote{برای فهرست تصاویر و جداول، از \lr{\textbackslash listoffigures} و \lr{\textbackslash listoftables} استفاده کنید.}

\subsection{متن اصلی}
\label{sec:first-book-mainmatter}

متن اصلی با دستور \cmd{mainmatter} شروع می‌شود:

\begin{latin}
\begin{minted}{latex}
\mainmatter
\mainmattergeometry
\end{minted}
\end{latin}

دستور \cmd{mainmatter} شماره‌گذاری صفحات را به
اعداد عربی (۱، ۲، ۳، ...) تغییر می‌دهد.
دستور \cmd{mainmattergeometry} نیز هندسه‌ی
صفحه را به حالت حاشیه‌ی پهن برمی‌گرداند.

\subsection{فصل}
\label{sec:first-book-chapter}

فصل با دستور \cmd{chapter} شروع می‌شود:

\begin{latin}
\begin{minted}{latex}
\chapter{|\pc{فصل اول}|}
\label{chap:first}
\end{minted}
\end{latin}

\mnote{دستور \lr{\textbackslash label} برای ارجاع به فصل استفاده می‌شود. همیشه یک برچسب معنادار بدهید.}

\subsection{پس‌متن}
\label{sec:first-book-backmatter}

پس‌متن با دستور \cmd{backmatter} شروع می‌شود:

\begin{latin}
\begin{minted}{latex}
\backmatter
\frontmattergeometry
\end{minted}
\end{latin}

در پس‌متن، معمولاً منابع و نمایه قرار می‌گیرند.

\subsection{جلد عقب}
\label{sec:first-book-backcover}

جلد عقب با دستور \cmd{makebackcover} ساخته
می‌شود:

\begin{latin}
\begin{minted}{latex}
\makebackcover{|\pc{توضیحات پشت جلد...}|}
\end{minted}
\end{latin}

\mnote{دستور \lr{\textbackslash makebackcover} باید حتماً در انتهای سند باشد، پیش از \lr{\textbackslash end\{document\}}.}

\section{افزودن منابع و نمایه}
\label{sec:first-book-bibliography-index}

کتاب شما اکنون کامل است، ولی هنوز منابع و نمایه
ندارد. در این بخش، آن‌ها را اضافه می‌کنیم.

\subsection{افزودن کتاب‌نامه}
\label{sec:first-book-bibliography}

برای افزودن کتاب‌نامه، در بخش پیش از
\cmd{begin\{document\}}، این خط را اضافه کنید:

\begin{latin}
\begin{minted}{latex}
\addbibresource{references.bib}
\end{minted}
\end{latin}

و در پس‌متن، این دستور را اضافه کنید:

\begin{latin}
\begin{minted}{latex}
\printbibliography[title={|\pc{منابع و مراجع}|}]
\end{minted}
\end{latin}

\mnote{کتاب‌نامه را داخل \lr{latin} قرار ندهید. \lr{xepersian} جهت متن را به‌درستی مدیریت می‌کند.}

\subsection{افزودن نمایه}
\label{sec:first-book-index}

برای افزودن نمایه، کافی است یک کلمه را با دستور
\cmd{index} علامت‌گذاری کنید:

\begin{latin}
\begin{minted}{latex}
|\pc{برنامه‌نویسی}|\index{|\pc{برنامه‌نویسی}|}
\end{minted}
\end{latin}

\mnote{دستور \lr{\textbackslash index} متن فارسی را مستقیماً در متن عادی می‌پذیرد و نیازی به \lr{\textbackslash pc} در متن عادی ندارد.}

و در پس‌متن، این دستور را اضافه کنید:

\begin{latin}
\begin{minted}{latex}
\printindex
\end{minted}
\end{latin}

\mnote{نمایه با \lr{xindy} ساخته می‌شود و مرتب‌سازی فارسی را به‌درستی انجام می‌دهد.}

\section{ساختار فایل‌های یک پروژه}
\label{sec:first-book-project-structure}

برای کتاب‌های بزرگ‌تر، بهتر است فایل‌ها را به
چند بخش تقسیم کنید:

\begin{latin}
\begin{minted}{text}
my-book/
├── main.tex
├── references.bib
├── chapters/
│   ├── chapter-01.tex
│   ├── chapter-02.tex
│   └── chapter-03.tex
├── figures/
│   ├── figure-01.pdf
│   └── figure-02.pdf
└── bibliography/
    └── references.bib
\end{minted}
\end{latin}

\mnote{تقسیم فایل‌ها به بخش‌های کوچک، مدیریت پروژه را بسیار آسان‌تر می‌کند. هر فصل در یک فایل جداگانه.}

سپس در \file{main.tex}:

\begin{latin}
\begin{minted}{latex}
\input{chapters/chapter-01}
\input{chapters/chapter-02}
\input{chapters/chapter-03}
\end{minted}
\end{latin}

\mnote{دستور \lr{\textbackslash input} فایل را در همان‌جا وارد می‌کند. برای فایل‌های بزرگ، از \lr{\textbackslash include} استفاده کنید که صفحه‌ی جدید ایجاد می‌کند.}

\section{خطاهای رایج}
\label{sec:first-book-mistakes}

\subsection{فراموش کردن \cmd{frontmatter} یا \cmd{mainmatter}}
\label{sec:first-book-mistake-frontmain}

اگر \cmd{frontmatter} و \cmd{mainmatter} را
فراموش کنید، شماره‌گذاری صفحات به‌درستی انجام
نمی‌شود و ممکن است شماره‌ها به‌صورت عربی از
ابتدای سند شروع شوند.

\textbf{راه‌حل:} ترتیب صحیح را رعایت کنید:
\cmd{frontmatter}، \cmd{mainmatter}، \cmd{backmatter}.

\subsection{قرار دادن \cmd{makebackcover} در ابتدای سند}
\label{sec:first-book-mistake-backcover}

اگر \cmd{makebackcover} را در ابتدای سند قرار
دهید، جلد عقب روی صفحه‌ی دوم ظاهر می‌شود و
بقیه‌ی کتاب به‌هم می‌ریزد.

\textbf{راه‌حل:} همیشه \cmd{makebackcover} را در
انتهای سند، پیش از \cmd{end\{document\}} قرار دهید.

\subsection{فراموش کردن \cmd{addbibresource}}
\label{sec:first-book-mistake-bibresource}

اگر \cmd{addbibresource} را فراموش کنید،
کتاب‌نامه خالی خواهد بود و استنادها به‌صورت
«؟؟» نمایش داده می‌شوند.

\textbf{راه‌حل:} همیشه \cmd{addbibresource} را
در پیش از \cmd{begin\{document\}} قرار دهید.

\section{تمرین‌ها}
\label{sec:first-book-exercises}

\begin{exercise}
    فایل \lr{my-first-book.tex} را بسازید و
    آن را کامپایل کنید. آیا فایل \lr{PDF} بدون
    خطا ساخته می‌شود؟
    \label{exc:first-book-compile}
\end{exercise}

\begin{exercise}
    یک فصل دوم اضافه کنید و در فهرست مطالب
    آن را ببینید. آیا شماره‌گذاری صفحات به‌درستی
    انجام می‌شود؟
    \label{exc:first-book-second-chapter}
\end{exercise}

\begin{exercise}
    یک کتاب‌نامه‌ی کوچک با دو منبع بسازید
    و با \cmd{cite} به آن‌ها استناد کنید.
    \label{exc:first-book-bibliography}
\end{exercise}

\begin{exercise}
    برای کتاب خود یک نمایه با سه مدخل بسازید.
    آیا مرتب‌سازی فارسی به‌درستی انجام می‌شود؟
    \label{exc:first-book-index}
\end{exercise}

\section{خلاصه}
\label{sec:first-book-summary}

در این فصل، اولین کتاب خود را ساختیم:

\begin{itemize}
    \item یک کتاب کامل در \lr{MatinBook} از پنج
    بخش تشکیل شده: جلد جلو، پیش‌متن، متن اصلی،
    پس‌متن، جلد عقب.

    \item ترتیب این بخش‌ها مهم است و باید رعایت
    شود.

    \item کتاب‌نامه با \cmd{addbibresource} و
    \cmd{printbibliography} اضافه می‌شود.

    \item نمایه با \cmd{index} و \cmd{printindex}
    اضافه می‌شود.

    \item برای کتاب‌های بزرگ، بهتر است فایل‌ها
    را به چند بخش تقسیم کنید.

    \item خطاهای رایج شامل فراموش کردن
    \cmd{frontmatter}، \cmd{mainmatter}،
    \cmd{backmatter}، \cmd{makebackcover} در
    انتهای سند، و \cmd{addbibresource} است.
\end{itemize}

در فصل بعدی (\cref{chap:book-structure})،
ساختار کامل یک کتاب را به‌طور مفصل بررسی
می‌کنیم و یاد می‌گیرید که هر بخش را چطور
شخصی‌سازی کنید.